\section{Introduction}\label{sec:intro}

In the Internet of Drones (IoD), drones send commands, telemetry and sensor data to ground stations over open radio links,
and ground stations forward data to cloud servers. An attacker can listen to, block, replay and modify these messages, and can
capture a drone and read its memory. Every session therefore needs mutual authentication, a fresh session key, and protection
of earlier and later sessions when a device is compromised.

BAKMM-IoD \citep{wazid2025} is a recent scheme for this setting. It provides drone--server and server--cloud authenticated key
agreement based on hashing and XOR, and a blockchain phase for the data. Its authors claim, among other properties, mutual
authentication, forward secrecy, untraceability and resistance to drone capture and stolen-verifier attacks, and support the
claims with an informal analysis and a Scyther verification.

\subsection{Motivation and contribution}

We re-examined BAKMM-IoD with two automated tools, Scyther \citep{cremers2008} and Tamarin \citep{meier2013}, and found that several of these claims do not hold in our Scyther and Tamarin models
(\cref{sec:bakmm-analysis}). Following the structure of earlier improvement papers such as \citet{li2024}, this paper reports
the findings and proposes and verifies an improved scheme. Our contributions are:
\begin{enumerate}
  \item We show that the published Scyther listing of BAKMM-IoD never reaches any of its claims, that the printed equations of
        the drone--server and server--cloud phases are inconsistent, and report these and further attacks as findings
        (\NumFlawsConfirmed\ confirmed, \NumFlawsConditional\ conditional), each with a tool result and an executable control model.
  \item We propose FSL-AKE v3 (\emph{v3} below; its predecessor in this project is called v2), a three-message scheme without timestamps that replaces XOR masking by authenticated
        encryption, separates every hash by a label, erases the old key at each commit and evolves it one-way, and repairs a
        lost message with a keyed one-time pseudonym.
  \item We verify v3 with Tamarin (\NumVthreeVerified\ of \NumVthreeLemmas\ drone--server lemmas verified, bounded to one
        session per pair and two for three lemmas, the
        \NumVthreeFalsified\ remaining ones being stated limits), with Scyther, and with a hand proof in the Real-Or-Random style,
        and we compare v3 with BAKMM-IoD lemma by lemma under the same attacker.
  \item We compare cost with the BAKMM-IoD paper's own method and timings and with a desktop benchmark of both reference
        implementations, and we report every quantity, including those where v3 is higher.
\end{enumerate}

\subsection{Threat model}\label{sec:threat}
The attacker controls the network, as in the Dolev--Yao (DY) model \citep{dolev1983} that BAKMM-IoD uses: it reads, blocks,
replays, reorders and injects messages. In addition, as in the Canetti--Krawczyk (CK) adversary model \citep{canetti2001} that
BAKMM-IoD adopts (its Sec.~3.2), it may reveal, at any time unless a
property restricts it: a session key; the ephemeral values of a session; the current memory of a drone (device capture); the
server database; the server's volatile state; and the server's long-term secret. An insider at the registration authority
(RA) and a leak of the server database are covered by these reveals. Denial of service, side channels, timing, and the
consensus of the blockchain phase are out of scope. The blockchain phase and handover between ground stations are outside the
scheme we propose.

\subsection{Organisation}
\Cref{sec:prelim} gives notation and tools. \Cref{sec:review} reviews BAKMM-IoD and \cref{sec:bakmm-analysis} analyses it.
\Cref{sec:scheme} presents FSL-AKE v3, \cref{sec:security} its security analysis and \cref{sec:performance} its performance.
\Cref{sec:discussion} discusses limitations and \cref{sec:conclusion} concludes.
